\hypertarget{appendix-a65290d66ed9}{}
\section{Shift-and-add multiplication}\label{appendix-a65290d66ed9-shift-and-add-multiplication}

\textbf{Curriculum depth:} Introduced · \textbf{First seen:} C002 · \textbf{Latest development:} C002

\subsection{Definition and mechanism}\label{appendix-a65290d66ed9-definition-and-mechanism}

Each set multiplier bit selects a shifted multiplicand partial product; an exact product of two \(N\)-bit values may need \(2N\) bits.

Multiplier bit i selects multiplicand shifted i positions. 13×11 uses 11=\texttt{1011₂}, hence 13+26+104=143. Keep the shifting multiplicand in the widened type; shifting a uint16\_t variable repeatedly discards high partial-product bits even if the accumulator is uint32\_t.

\subsection{Prerequisites}\label{appendix-a65290d66ed9-prerequisites}

\begin{itemize}
\tightlist
\item
  \hyperref[appendix-73425b25118e]{Shift operations} --- Logical right shift fills high positions with zero; arithmetic right shift replicates the sign bit.
\item
  \hyperref[appendix-d1a5b1e520c2]{Full adder} --- A full adder adds A, B, and carry-in.
\end{itemize}

\subsection{Used by}\label{appendix-a65290d66ed9-used-by}

Connect later chunks here when they use this concept.

\subsection{Recall and next depth}\label{appendix-a65290d66ed9-recall-and-next-depth}

Explain shift-and-add multiplication without the label, work the example above, then state the assumption whose removal would change the result.

\textbf{Pending:} Booth recoding, compressor trees, pipelining, and signed multiplication.

\subsection{Source and navigation}\label{appendix-a65290d66ed9-source-and-navigation}

\hyperref[01-foundations-i-ch001-arithmetic-becomes-a-policy]{C002 teaching chapter} · \hyperref[appendix-fd090d982a15]{Part definitions and mechanisms} · \href{../02-Foundations-II/main.pdf\#nameddest=appendix-4ad47bdd269a}{Dependency Map (02-Foundations-II, appendix-4ad47bdd269a)}
